\chapter{2017 Nobel Prize in Economics}
\textbf{Laureates: }Richard Thaler (USA)

\section*{Reason for Award}
For his contributions to behavioral economics

\section{What They Did}
Richard Thaler advanced behavioral economics by incorporating psychological
insights into economic analysis. Born in East Orange, New Jersey, in 1945,
Thaler studied at Case Western Reserve University and the University of Rochester
before becoming a professor at the University of Chicago Booth School of Business.

He identified systematic deviations from rational behavior, including mental
accounting, where people treat money differently depending on its source or
intended use. Thaler showed that people maintain separate mental accounts for
different purposes, leading to decisions that depart from the standard assumption
that money is fully fungible. People often budget by category rather than treating
all income as interchangeable.

Thaler documented the endowment effect, showing that people often value items
more highly when they own them. In experiments, subjects given a mug typically
demanded more to sell it than they were willing to pay to buy one, a pattern
consistent with loss aversion and reference-dependent preferences. This finding
challenged the standard assumption that preferences are independent of endowment.

He analyzed present bias, demonstrating how people discount future rewards more
steeply than the standard exponential model predicts. Thaler showed that people
often prefer immediate gratification over larger future rewards, leading to
procrastination, inadequate saving, and unhealthy behaviors. Later formal
quasi-hyperbolic models built on this behavioral insight and have been used in
economics and public policy.

Thaler's concept of nudging showed how choice architecture can help people
make better decisions without restricting freedom. He argued that since choice
architecture is unavoidable, it should be designed to help people achieve their
own goals. His book ``Nudge'' (2008), co-authored with Cass Sunstein, popularized
these ideas and influenced policy design worldwide.

Thaler also studied social preferences, including fairness and reciprocity. His
experimental and theoretical work helped explain why people may reject unfair
offers or resist price increases that they consider unjust, even when doing so
is costly to themselves.

His book ``Misbehaving'' (2015) traced the development of behavioral economics
as a field, documenting the resistance he faced from mainstream economists
who viewed his work as anathema to the rational agent paradigm. Thaler's
advocacy helped establish behavioral economics as a legitimate and important
field of study.

\section{Impact on Economic Thinking}
Thaler transformed economics by showing that traditional assumptions about
rationality can be too restrictive for some questions. His work demonstrated
that people may have limited self-control, limited forethought, and social
preferences that are not purely self-interested. These insights helped explain
puzzles such as the equity-premium puzzle, the disposition effect in trading,
and persistent differences in saving behavior.

The endowment effect, mental accounting, and present bias have been examined in
many experiments and empirical studies, establishing behavioral economics as a
substantial research program rather than a theoretical curiosity. These findings
have influenced finance, marketing, public policy, and organizational behavior.

Thaler's nudging concept influenced policy design, including automatic enrollment
in retirement plans, organ-donor registration, and efforts to encourage healthier
choices. The UK's Behavioral Insights Team, established in 2010, became an
important example of a government unit applying behavioral insights to policy
challenges.

His work bridged economics and psychology, creating a more psychologically
realistic foundation for understanding human behavior. The recognition of bounded
rationality encouraged policy interventions that take people's actual decision-
making processes into account rather than assuming perfect rationality.

Thaler's contributions earned him recognition as a pioneer who made economics
more relevant to real-world problems. He demonstrated that economics can be
both scientifically rigorous and practically useful, opening the door for
interdisciplinary research that combines insights from psychology, neuroscience,
and computer science.

\section{Modern Perspectives Today}
Behavioral economics has become mainstream, with Thaler's insights influencing
policy, business, and academic research. Automatic-enrollment programs have often
increased participation in retirement savings, although their effects depend on
the contribution rates, defaults, and populations involved.

Nudge units in governments around the world apply behavioral insights to policy
design. International development organizations have also incorporated behavioral
economics into some programs. Financial products increasingly use behavioral
devices such as default options, simplified disclosures, and automatic escalation
of contributions.

The field has expanded to address climate change, public health, and education
policy. Behavioral researchers also study how well-being and flourishing can
inform policy beyond traditional economic measures. Thaler's concept of ``choice
architecture'' has influenced product design, website interfaces, and
organizational decision-making processes.

As behavioral economics matures, its insights continue to reshape how economists
understand decision-making, design institutions, and evaluate policy interventions.
The 2017 Nobel Prize recognized Thaler's transformative contribution to making
economics more realistic and relevant to the world as it actually is.

\subsection*{Key References}
\begin{itemize}
    \item Thaler, R. H. (1980). ``Toward a Positive Theory of Consumer Choice.'' \textit{Journal of Economic Behavior \& Organization}, 1(1), 39--60.
    \item Thaler, R. H., \& Sunstein, C. R. (2008). \textit{Nudge: Improving Decisions About Health, Wealth, and Happiness}. Yale University Press.
\end{itemize}
